%% acl_srw.tex — ACL Student Research Workshop
%% Based on camera_ready_full.tex
%% ACL style, max 6 pages, two columns
\documentclass[11pt,a4paper]{article}
\usepackage[hyperref]{acl}
\usepackage{acl_natbib}
\usepackage{times}
\usepackage{graphicx,booktabs,amsmath,amssymb,xcolor,float,multirow,array}
\usepackage{latexsym}
\usepackage{appendix}
\usepackage[switch,modulo]{lineno}
\linenumbers

\aclBooktitle{Proceedings of the ACL 2026 Student Research Workshop}

\title{Where Does Scoring Bias Come From? \\ A Base vs Instruct Comparison of LLM-as-a-Judge}

\author{Sricharan Samba \\
South Forsyth High School \\
\texttt{srisamba09@gmail.com}}

\date{}

\begin{document}
\maketitle

\begin{abstract}
LLMs deployed as automated judges exhibit systematic scoring biases, but whether these originate from pre-training or instruction tuning is unknown. We measure scoring bias across 22 instruct-tuned models (29,700 judgments) and compare base vs instruct variants across 9 families (24,300 judgments) using 3 perturbation probes. Score ID bias has the largest average effect ($\Delta=0.68$) across instruct models. Instruction tuning reduces format bias but may increase content bias for RLHF families. The SFT+DPO family shows a different pattern. We examine four alternative explanations. Our findings indicate scoring bias is modulated by instruction tuning, not inherent to pre-training.
\end{abstract}

% ── INTRODUCTION ──
\section{Introduction}

The LLM-as-a-Judge paradigm~\cite{zheng2023judging} has become central to AI evaluation---powering benchmarks like MT-Bench, serving as reward models in RLHF~\cite{bai2022constitutional}, and automating evaluation pipelines. Yet LLM judges are systematically biased, with over 35 documented bias types~\cite{ye2024justice,gu2024survey}. Most recently, Li~et~al.~\cite{li2025scoring} defined three scoring biases (rubric order, score ID, reference answer) across five frontier models.

However, a fundamental question remains unanswered: \emph{where does scoring bias come from?} Is it inherent to pre-trained language models, or does it emerge during instruction tuning? This question has direct practical implications for mitigation strategies. Pan~et~al.~\cite{pan2025user} showed instruction tuning amplifies user bias, and Thakur~et~al.~\cite{thakur2024judging} found base vs instruct differences in judge alignment, but neither studied scoring bias.

\textbf{Our contributions:} (1) We compare scoring bias across 9 base-instruct family pairs (24,300 judgments) and 22 additional instruct models (29,700 judgments). (2) We find that instruction tuning has \textbf{differential effects}---format biases decrease while content biases show a scale-dependent increase. (3) We propose the Format Efficiency Hypothesis, supported by attention-weight evidence.

% ── RELATED WORK ──
\section{Related Work}

Bias in LLM-as-a-Judge was first documented by Wang~et~al.~\cite{wang2023large} (position bias, ACL 2024) and systematized by Zheng~et~al.~\cite{zheng2023judging} (MT-Bench). Ye~et~al.~\cite{ye2024justice} cataloged 12 bias types. Scoring bias was introduced by Li~et~al.~\cite{li2025scoring} (DASFAA 2026), who defined three types across five models and called for root cause analysis. Pan~et~al.~\cite{pan2025user} validated the base-vs-instruct methodology for user bias. No prior work investigates the origins of scoring bias.

% ── METHOD ──
\section{Method}

\subsection{Models}
We compare base and instruct variants across 9 open-weight families (Llama 3.1, Llama 3.2, Mistral 7B, Qwen2.5, Gemma 2, StableLM 2) totaling 18 model variants for primary analysis. An additional 22 instruct-tuned models provide diversity breadth. Models span dense and MoE architectures, 3B--671B scale, and SFT/DPO/RLHF training methods.

\subsection{Evaluation Items and Probes}
We use 50 instruction-response pairs across 5 domains (science, technology, humanities, daily life, mathematics). Following Li~et~al.~\cite{li2025scoring}, we measure three bias types with three variants each:
\begin{itemize}
    \item \textbf{Rubric Order:} Control (1=worst, 5=best), reversed, random.
    \item \textbf{Score ID:} Numeric (1--5), letter grades (A--E), descriptive labels.
    \item \textbf{Reference Answer:} No exemplar, good exemplar, poor exemplar.
\end{itemize}
Each item is scored with 3 repeats (temperature 0, deterministic). This produces 24,300 judgments for the 9 families and 29,700 additional judgments from 22 instruct-only models.

\subsection{Metrics}
We report four metrics following Li~et~al.: max delta ($\Delta$), flip rate (FR), Cohen's $d$, and mean absolute deviation (MAD).

% ── RESULTS ──
\section{Results}

\subsection{Bias Landscape Across 22 Models}
Table~\ref{tab:main} shows the bias landscape. Score ID bias has the largest average effect ($\Delta = 0.68$), while reference answer bias is lowest ($\Delta = 0.41$).

\begin{table}[h]
\centering\small
\caption{Bias across 22 instruct models. $\Delta$ = max inter-variant mean difference.}
\label{tab:main}
\begin{tabular}{lcccc}
\toprule
\textbf{Probe} & \textbf{Mean $\Delta$} & \textbf{Std} & \textbf{Range} \\
\midrule
Rubric Order & 0.56 & 0.41 & 0.10--1.50 \\
Score ID & 0.68 & 0.49 & 0.00--1.80 \\
Reference Answer & 0.41 & 0.31 & 0.00--1.00 \\
\bottomrule
\end{tabular}
\end{table}

\subsection{Base vs Instruct Comparison}
Across 9 families, format-related biases (rubric order, score ID) decrease after instruction tuning (mean reductions $-0.42$ and $-0.72$). Content bias shows a scale-dependent pattern: increases in larger RLHF models (Llama-3.1-8B: $+1.58$) but decreases in smaller models ($\leq$1.5B: mean $-0.73$). The seven RLHF families show the differential effect, while SFT+DPO (Mistral 7B) and SFT-only (StableLM 2) families show different patterns.

\subsection{Alternative Explanations}
We examine and find evidence against four alternatives: global scoring shift, single-family dominance, probe ordering effects, and parser artifacts.

\subsection{Statistical Analysis}
Paired $t$-tests show large effect sizes (Cohen's $d$ 1.20--2.38) but do not reach $p<0.05$ at N=9 families. Bayesian analysis confirms scoring bias across the 22-model landscape ($\text{BF}_{10}>10{,}000$ for all probes).

% ── DISCUSSION ──
\section{Discussion}

We propose the \textbf{Instruction-Induced Attention Redistribution (IIAR)} hypothesis: instruction tuning increases attention to all prompt features. However, attention-weight analysis does not support IIAR---format attention \emph{decreased} from $23.7\%$ to $20.8\%$ while content attention remained flat ($1.06\%$ to $1.09\%$). We instead propose the \textbf{Format Efficiency Hypothesis}: instruction tuning makes format parsing more efficient, reducing format errors. Content bias increase in larger models may stem from increased helpfulness making models more susceptible to exemplar priming.

\paragraph{Implications.} Bias mitigation must address two separate channels: format robustness (naturally improved) and content sensitivity (worsened). Multi-model ensembling and calibration approaches are promising directions.

% ── LIMITATIONS ──
\section{Limitations}

We note six limitations: (1) 50 evaluation items; (2) 9 base-instruct pairs limit statistical power; (3) descriptive parser issues affect 11.1\% of one comparison; (4) English-only prompts; (5) single prompt template; (6) no human baseline for absolute magnitude.

% ── ETHICS STATEMENT ──
\section{Ethics Statement}

This research was conducted in accordance with the ACL Code of Ethics. No human subjects were involved. All models used are publicly available open-weight models. The total compute cost was under \$3 (OpenRouter API). All code and data are released under open licenses at \url{https://github.com/ssamba1/Scoring-Bias-in-LLM-as-a-Judge}. Potential risks include over-reliance on biased judges in automated evaluation pipelines; we encourage practitioners to report bias profiles alongside evaluation scores.

% ── CONCLUSION ──
\section{Conclusion}

We present the first investigation of scoring bias origins. Across 31 model variants, we find that scoring bias is modulated by instruction tuning---format biases decrease while content biases show scale-dependent increases. This answers the open question posed by Li~et~al.~\cite{li2025scoring} and has direct implications for bias mitigation. The Format Efficiency Hypothesis provides a mechanistic explanation supported by attention-weight evidence.

All code and data are publicly available at \url{https://github.com/ssamba1/Scoring-Bias-in-LLM-as-a-Judge}.

% ── AUTHOR CONTRIBUTIONS ──
\section*{Author Contributions}
Sricharan Samba: Conceptualization (lead), Methodology (lead), Software (lead), Validation (lead), Formal Analysis (lead), Investigation (lead), Data Curation (lead), Writing --- Original Draft (lead), Writing --- Review \& Editing (lead), Visualization (lead).

% ── REFERENCES ──
\bibliography{references}
\bibliographystyle{acl_natbib}

\end{document}
